% ---- pg-277 （书页 265） ----

\begin{mvtable}
{\zihao{6}\hfill 续表\par}\vspace{2pt}
\begin{tabular}{ccc|ccc}
\toprule
\multicolumn{3}{c}{汽油卡车} & \multicolumn{3}{c}{柴油卡车} \\
\cmidrule(lr){1-3}\cmidrule(lr){4-6}
$x_{1}$ & $x_{2}$ & $x_{3}$ & $x_{1}$ & $x_{2}$ & $x_{3}$ \\
\midrule
10.25 & 5.07 & 10.17 & 9.09 & 13.25 & 20.66 \\
11.11 & 6.15 & 7.61 & 8.53 & 10.14 & 17.45 \\
12.17 & 14.26 & 14.39 & 8.29 & 6.22 & 16.38 \\
10.24 & 2.59 & 6.09 & 15.90 & 12.90 & 19.09 \\
10.18 & 6.05 & 12.14 & 11.94 & 5.69 & 14.77 \\
8.88 & 2.70 & 12.23 & 9.54 & 16.77 & 22.66 \\
12.34 & 7.73 & 11.68 & 10.43 & 17.65 & 10.66 \\
8.51 & 14.02 & 12.01 & 10.87 & 21.52 & 28.47 \\
26.16 & 17.44 & 16.89 & 7.13 & 13.22 & 19.44 \\
12.95 & 8.24 & 7.18 & 11.88 & 12.18 & 21.20 \\
16.93 & 13.37 & 17.59 & 12.03 & 9.22 & 23.09 \\
14.70 & 10.78 & 14.58 & & & \\
10.32 & 5.16 & 17.00 & & & \\
8.98 & 4.49 & 4.26 & & & \\
9.70 & 11.59 & 6.83 & & & \\
12.72 & 8.63 & 5.59 & & & \\
9.49 & 2.16 & 6.23 & & & \\
8.22 & 7.95 & 6.72 & & & \\
13.70 & 11.22 & 4.91 & & & \\
8.21 & 9.85 & 8.17 & & & \\
15.86 & 11.42 & 13.06 & & & \\
9.18 & 9.18 & 9.49 & & & \\
12.49 & 4.67 & 11.94 & & & \\
17.32 & 6.86 & 4.44 & & & \\
\bottomrule
\end{tabular}
\end{mvtable}

{\zihao{6} 资料来源：数据经 M. Keaton 许可使用.}

(a) 取 $\alpha=0.01$, 试检验两个费用均值向量是否相等.

(b) 如果 (a) 中的费用均值向量的相等性假设被拒绝, 试找出导致拒绝该假设的均值分量的线性组合.

(c) 对每对均值分量, 构造其 99% 联合置信区间, 哪些费用看上去有显著差异(若有的话)?

(d) 对你在分析中所用的那些假定的正确性作出评论. 特别应注意汽油卡车的第 9 和第 21 个观测值被识别为多元离群值(参见练习 5.22 和文献[2]). 将这两观测值删除, 重做 (a) 的检验. 对所得结果作出评论.

\medskip
\noindent 6.20\quad 对 45 只雄性钩镰鸢的尾长 $x_{1}$(毫米)和翅长 $x_{2}$(毫米)进行测量, 其结果见\hyperref[tab:6.11]{表 6.11}, 对雌性钩镰鸢的尾长及翅长的测量结果已在\hyperref[tab:5.12]{表 5.12} 中给出.

(a) 画出雄性钩镰鸢数据的散布图, 并检验有无离群数据. (特别注意第 31 个观察值 $x_{1}=284$).

(b) 检验雄性和雌性钩镰鸢两个总体均值向量的相等性(取 $\alpha=0.05$). 如果 $H_{0}: \boldsymbol{\mu}_{1}-\boldsymbol{\mu}_{2}=\mathbf{0}$ 被拒绝, 试找出导致拒绝 $H_{0}$ 的均值分量的线性组合. (对雄性钩镰鸢数据, 在进行检验之前
